\chapter{Conclusion}
\label{ch:Conclusion}

\noindent This project developed a computational implementation of H\"uckel molecular orbital theory to investigate the electronic structure of finite graphene flakes. The model successfully reproduced the analytical energy spectrum of benzene and demonstrated close agreement with published H\"uckel calculations for larger graphene structures. For completed shell flakes, the HOMO-LUMO gap decreased systematically from 4.80 eV for benzene ($N=6$) to approximately 0.35 eV for $N=294$, a reduction of 93\%. A power-law fit yielded $\Delta E_{\mathrm{HL}}\propto N^{-0.660}$, indicating a strong size dependence, although the fitted exponent differs from the idealised $N^{-1/2}$ scaling estimate. \\

\noindent The investigation of incomplete outer shells revealed a more complex, non-monotonic dependence on carbon site count, demonstrating that electronic properties are influenced by lattice connectivity and symmetry rather than size alone. Notably, the $N=32$ structure exhibited two zero-energy states associated with a sublattice imbalance of $|N_A-N_B|=2$. Furthermore, the reconstruction of molecular orbitals using Slater-type functions provided a qualitative visualisation of their spatial distributions, phases and bonding characteristics, complementing the numerical energy analysis. \\

\noindent Despite its simplicity, the nearest neighbour H\"uckel approximation provides a useful framework for understanding electronic trends in finite conjugated carbon systems. Its limitations, particularly the neglect of electron-electron interactions, structural relaxation and realistic edge termination, restrict direct comparisons with experimental electronic properties. Future extensions could incorporate these effects or compare the present results with higher level electronic structure calculations. Nevertheless, the project demonstrates how a relatively simple computational model can reveal the relationship between atomic structure and electronic behaviour in finite graphene flakes.
